%% 2-18 表計算で計算しよう（第11回：ページ制限なし。なぜ番地がずれるのか、なぜその数なのかを1段ずつ）
%% この章の部品（頭文字 \bkiis…）
\makeatletter
% 式を入れる欄（セルの名前と、入れた式）  \bkiisbar{番地}{式}
\newcommand\bkiisbar[2]{\tikz[baseline=(t.base), x=1mm, y=1mm]{%
  \draw[bkGray, line width=0.6pt, fill=white] (0,-2.6) rectangle (12,3.4);
  \node[font=\sffamily\small, text=bkInk] (t) at (6,0) {#1};
  \node[font=\sffamily\itshape\footnotesize, text=bkGray] at (16,0.2) {fx};
  \draw[bkGray, line width=0.6pt, fill=white] (20,-2.6) rectangle (62,3.4);
  \node[anchor=west, font=\ttfamily\small, text=bkInk] at (22,0.2) {#2};}}
% 選んだセルの色
\colorlet{bkiissel}{bkYellow!35}
\makeatother

\bkchapter{2}{18}{\R{表計算}{ひょう|けい|さん}で\R{計算}{けい|さん}しよう}{\bklevel{2}}{\R{表計算}{ひょう|けい|さん}で\R{計算}{けい|さん}しよう}
\bkmichishirube{1-18,2-11,2-13}{5-8,6-2}

\begin{goal}
\begin{itemize}[label={\color{bkTheme}●}, leftmargin=1.6\zw, itemsep=0pt, topsep=1pt]
\item \bklevel{2}\enspace セルの\R{番地}{ばん|ち}を\R{使}{つか}って「＝」で\R{始}{はじ}まる\R{式}{しき}を\R{書}{か}き、\R{式}{しき}をコピーして\R{同}{おな}じ\R{計算}{けい|さん}をくり\R{返}{かえ}せる。
\item \bklevel{2}\enspace ずらしたくない\R{番地}{ばん|ち}に \$ をつけて\R{使}{つか}える。
\item \bklevel{2}\enspace \R{上}{うえ}の\R{行}{ぎょう}の\R{結果}{けっ|か}を\R{使}{つか}って\R{下}{した}の\R{行}{ぎょう}を\R{計算}{けい|さん}し、\R{表}{ひょう}とグラフで\R{変}{か}わり\R{方}{かた}を\R{調}{しら}べられる。
\end{itemize}
\end{goal}

\begin{hatena}
\noindent\begin{minipage}[c]{0.5\linewidth}
お\R{年玉}{とし|だま}の3000\R{円}{えん}を\R{貯金箱}{ちょ|きん|ばこ}に\R{入}{い}れ、\R{毎月}{まい|つき}500\R{円}{えん}ずつ\R{入}{い}れていきます。ちょうど10000\R{円}{えん}になるのは、\R{何}{なん}か\R{月後}{げつ|ご}かな？　10\R{年後}{ねん|ご}には、いくらになるかな？
\end{minipage}\hfill
\begin{minipage}[c]{0.47\linewidth}\centering
\begin{tikzpicture}[x=1mm, y=1mm]
  \foreach \i/\n/\lab in {0/6/0, 1/7/1, 2/8/2, 3/9/3} {
    \foreach \k in {1,...,\n} {\fill[bkYellow!85!black] (\i*10,\k*2.2) ellipse (3.6 and 1.1); \draw[bkYellow!50!black, line width=0.4pt] (\i*10,\k*2.2) ellipse (3.6 and 1.1);}
    \node[anchor=north, font=\sffamily\footnotesize, text=bkInk] at (\i*10,-1.2) {\lab};}
  \node[font=\sffamily\bfseries\large, text=bkGray] at (38.5,9) {…};
  \node[font=\sffamily\bfseries\Large, text=bkTheme] at (46,11) {？};
  \node[anchor=north, font=\sffamily\footnotesize, text=bkInk] at (46,-1.2) {120};
  \node[anchor=north, font=\sffamily\footnotesize, text=bkInk] at (21,-7) {（か\R{月}{げつ}）};
  \node[anchor=south west, font=\sffamily\footnotesize, text=bkInk, inner sep=0pt] at (-3.6,24) {500\R{円玉}{えん|だま}1\R{枚}{まい}が \tikz[baseline=-0.6ex]{\fill[bkYellow!85!black] (0,0) ellipse (1.8mm and 0.6mm);}};
\end{tikzpicture}
\end{minipage}

\vspace{4pt}
10\R{年}{ねん}は120か\R{月}{げつ}です。\R{電卓}{でん|たく}で1か\R{月}{げつ}ずつたしていくと、120\R{回}{かい}もたし\R{算}{ざん}をすることになります。
\end{hatena}

\noindent 1つ1つは\R{同}{おな}じ「500をたす」\R{計算}{けい|さん}です。\R{同}{おな}じ\R{計算}{けい|さん}を\R{何百回}{なん|びゃっ|かい}もくり\R{返}{かえ}すのは、\R{人}{ひと}には\R{大変}{たい|へん}ですが、コンピュータには\R{得意}{とく|い}なことです。この\R{章}{しょう}では、\textbf{\R{表計算}{ひょう|けい|さん}ソフト}を\R{使}{つか}って、くり\R{返}{かえ}しの\R{計算}{けい|さん}をコンピュータにまかせる\R{方法}{ほう|ほう}を\R{学}{まな}びます。\bkidx{ひょうけいさんそふと}{\R{表計算}{ひょう|けい|さん}ソフト}

%% ---------------- やってみよう・思い出そう ----------------
\begin{yattemiyou}[\R{表計算}{ひょう|けい|さん}ソフトに\R{計算}{けい|さん}させよう]
タブレットやパソコンで\R{表計算}{ひょう|けい|さん}ソフトを\R{開}{ひら}きます。\R{画面}{が|めん}は、たてと\R{横}{よこ}に\R{線}{せん}を\R{引}{ひ}いた\R{大}{おお}きな\R{表}{ひょう}です。
\begin{enumerate}[label={\bkcircled[bkGreen]{\arabic*}}, leftmargin=2.4\zw, itemsep=1pt, topsep=3pt]
\item A\R{列}{れつ}の1\R{行目}{ぎょう|め}（A1）に3、2\R{行目}{ぎょう|め}（A2）に5を\R{入}{い}れよう。
\item A3に「\texttt{=A1+A2}」と\R{入}{い}れて、Enterキーをおそう。A3には\R{何}{なに}が\R{出}{で}るかな？
\item A1を10に\R{変}{か}えよう。A3はどうなるかな？
\end{enumerate}

\begin{bkfigure}
\begin{tikzpicture}[x=1mm, y=1mm]
  \node[anchor=west, inner sep=0pt] at (0,30) {\bkiisbar{A3}{=A1+A2}};
  \node[anchor=north west, inner sep=0pt] at (0,24) {\bksheet{1}{18}{%
1 & 3 \\
2 & 5 \\
3 & \SetCell{bg=bkiissel} 8 \\}};
  \node[font=\sffamily\small, text=bkInk, align=center] at (48,14) {A1を10に\\\R{変}{か}えると};
  \draw[bk arrow, bkGreen] (38,5) -- (58,5);
  \node[anchor=north west, inner sep=0pt] at (66,24) {\bksheet{1}{18}{%
1 & \SetCell{bg=bkiissel} 10 \\
2 & 5 \\
3 & \SetCell{bg=bkiissel} 15 \\}};
\end{tikzpicture}
\bkfigcaption{A3には\R{式}{しき}の\R{答}{こた}えが\R{出}{で}る。A1を\R{変}{か}えると、A3の\R{答}{こた}えも\R{自動}{じ|どう}で\R{変}{か}わる。}
\end{bkfigure}
\end{yattemiyou}

\noindent A3には $3+5$ の\R{答}{こた}え8が\R{出}{で}ます。A1を10に\R{変}{か}えると、A3は\R{何}{なに}もしなくても $10+5=15$ に\R{変}{か}わります。A3には「8」という\R{数}{かず}ではなく、「A1とA2をたす」という\R{計算}{けい|さん}のしかたが\R{入}{はい}っているからです。\R{画面}{が|めん}の\R{上}{うえ}の\R{細長}{ほそ|なが}い\R{欄}{らん}（\R{式}{しき}の\R{欄}{らん}）を\R{見}{み}ると、A3に\R{入}{はい}っている\R{式}{しき}が\R{見}{み}えます。

\begin{omoidasou}
\bkrecall{1-18}{おれ\R{線}{せん}グラフにすると、\R{変}{か}わり\R{方}{かた}がひと\R{目}{め}でわかる。}
\bkrecall{2-11}{\R{計算}{けい|さん}の\R{順番}{じゅん|ばん}：かっこの\R{中}{なか}が\R{先}{さき}。×÷は＋－より\R{先}{さき}。$2+3\times4=14$、$(2+3)\times4=20$}
\bkrecall{2-13}{\R{式}{しき}の\R{文字}{も|じ}に\R{数}{かず}を\R{入}{い}れると、\R{式}{しき}の\R{値}{あたい}が\R{出}{で}る（\R{代入}{だい|にゅう}）。}
\bkrecall{2-5}{\qty{5}{\percent}は0.05、\qty{10}{\percent}は0.1。もとにする\R{量}{りょう}の\qty{5}{\percent}は、もとにする\R{量}{りょう}×0.05。}
\bkrecall{2-10}{1\R{時間}{じ|かん}あたり60\R{分}{ぷん}で \qty{60}{min/h}。1\R{時間}{じ|かん}あたり3600\R{秒}{びょう}で \qty{3600}{s/h}。1 kmあたり1000 mで \qty{1000}{m/km}。}
\end{omoidasou}

%% ---------------- 数学の道具 ----------------
\bksection{doguu}{\R{数学}{すう|がく}の\R{道具}{どう|ぐ}}

\bksubsection{1}{セルと\R{番地}{ばん|ち}、「＝」で\R{始}{はじ}まる\R{式}{しき}}

\leavevmode\label{bk:2-18:shiki}\noindent\begin{minipage}[c]{0.52\linewidth}
\R{表計算}{ひょう|けい|さん}の\R{表}{ひょう}の、1つ1つのますを\textbf{セル}といいます。\bkidx{せる}{セル}たての\R{列}{れつ}にはA・B・C…、\R{横}{よこ}の\R{行}{ぎょう}には1・2・3…と\R{名前}{な|まえ}がついています。B\R{列}{れつ}の3\R{行目}{ぎょう|め}のセルは「B3」です。これをセルの\textbf{\R{番地}{ばん|ち}}といいます。\bkidx{ばんち}{\R{番地}{ばん|ち}}\R{列}{れつ}の\R{文字}{も|じ}が\R{先}{さき}、\R{行}{ぎょう}の\R{番号}{ばん|ごう}があとです。\R{家}{いえ}の\R{住所}{じゅう|しょ}のように、\R{番地}{ばん|ち}がわかれば、どのセルかが1つに\R{決}{き}まります。
\end{minipage}\hfill
\begin{minipage}[c]{0.45\linewidth}\raggedleft
\begin{tikzpicture}[x=1mm, y=1mm]
  % 列の名前と行の番号
  \fill[bkLine!45] (0,0) rectangle (43,-5);
  \fill[bkLine!45] (0,0) rectangle (7,-29);
  \foreach \c [count=\k] in {A,B,C} \node[font=\sffamily\scriptsize] at (7+\k*12-6,-2.5) {\c};
  \foreach \r in {1,...,4} \node[font=\sffamily\scriptsize] at (3.5,-5-\r*6+3) {\r};
  \fill[bkiissel] (19,-17) rectangle (31,-23);
  \draw[bkTheme, line width=1.2pt] (19,-17) rectangle (31,-23);
  \foreach \x in {0,7,19,31,43} \draw[bkLine, line width=0.5pt] (\x,0) -- (\x,-29);
  \foreach \y in {0,-5,-11,-17,-23,-29} \draw[bkLine, line width=0.5pt] (0,\y) -- (43,\y);
  \node[font=\sffamily\bfseries\small, text=bkInk] at (25,-20) {B3};
  \draw[bkTheme, line width=0.9pt, -{Stealth[length=4pt]}] (25,4.8) -- (25,1.5);
  \node[anchor=south, font=\sffamily\footnotesize, text=bkTheme!85!black] at (25,5.3) {B\R{列}{れつ}};
  \draw[bkTheme, line width=0.9pt, -{Stealth[length=4pt]}] (-5.5,-20) -- (-1.5,-20);
  \node[anchor=east, font=\sffamily\footnotesize, text=bkTheme!85!black] at (-6,-20) {3\R{行目}{ぎょう|め}};
\end{tikzpicture}
\end{minipage}

\par\addvspace{6pt}
セルに「＝」から\R{始}{はじ}まる\R{式}{しき}を\R{入}{い}れると、\R{表計算}{ひょう|けい|さん}ソフトが\R{計算}{けい|さん}して、\R{答}{こた}えを\R{表示}{ひょう|じ}します。「＝」で\R{始}{はじ}めるのは、「これは\R{文字}{も|じ}ではなく\R{計算}{けい|さん}の\R{式}{しき}です」とコンピュータに\R{知}{し}らせるためです。

キーボードには×や÷のキーがありません。また、×はアルファベットのxとまちがえやすいので、\R{次}{つぎ}の\R{記号}{き|ごう}を\R{使}{つか}います。

\begin{bkcenter}
\begin{bktable}{colspec={Q[l,m,wd=24mm] *{5}{Q[c,m,wd=15mm]}}, rowsep=1.5pt, cells={font=\sffamily\small}, column{2-6}={font=\ttfamily\small}}
\R{算数}{さん|すう}の\R{記号}{き|ごう} & {\sffamily ＋} & {\sffamily －} & {\sffamily ×} & {\sffamily ÷} & {\sffamily $2^3$} \\
\R{表計算}{ひょう|けい|さん}の\R{記号}{き|ごう} & + & - & * & / & 2\^{}3 \\
\end{bktable}
\bktabcaption{\R{表計算}{ひょう|けい|さん}で\R{使}{つか}う\R{計算}{けい|さん}の\R{記号}{き|ごう}。/ は\R{分数}{ぶん|すう}の\R{横線}{よこ|せん}をななめにした\R{形}{かたち}。}
\end{bkcenter}

\R{計算}{けい|さん}の\R{順番}{じゅん|ばん}は、\R{算数}{さん|すう}と\R{同}{おな}じです（2-11）。\texttt{=2+3*4} と\R{入}{い}れると、*を\R{先}{さき}に\R{計算}{けい|さん}して $2+12=14$。\texttt{=(2+3)*4} と\R{入}{い}れると、かっこの\R{中}{なか}を\R{先}{さき}に\R{計算}{けい|さん}して $5\times4=20$ です。\R{表計算}{ひょう|けい|さん}の\R{式}{しき}も、\R{必要}{ひつ|よう}なところにはかっこをつけます。

\R{式}{しき}の\R{中}{なか}には、\R{数}{かず}のかわりにセルの\R{番地}{ばん|ち}を\R{書}{か}けます。たとえば、A2に\R{長方形}{ちょう|ほう|けい}のたて3、B2によこ4を\R{入}{い}れ、C2に \texttt{=A2*B2} と\R{入}{い}れると、C2に\R{面積}{めん|せき}12が\R{出}{で}ます（\R{次}{つぎ}のページの\R{図}{ず}の2\R{行目}{ぎょう|め}）。\R{番地}{ばん|ち}の\R{所}{ところ}に、そのセルの\R{数}{かず}が\R{代入}{だい|にゅう}されるのです（2-13）。\R{文字}{も|じ}の\R{式}{しき} $x\times y$ の $x$ と $y$ に\R{数}{かず}を\R{入}{い}れるのと\R{同}{おな}じです。

\bksubsection{2}{\R{式}{しき}をコピーすると、\R{番地}{ばん|ち}もずれる}

\leavevmode\label{bk:2-18:copy}\R{長方形}{ちょう|ほう|けい}が4つあるときは、C2の\R{式}{しき}をC3〜C5にコピーします（セルの\R{右下}{みぎ|した}の\R{角}{かど}を\R{下}{した}へひっぱる）。コピーした\R{式}{しき}はどうなるかな？

\begin{bkfigure}
\begin{tikzpicture}[x=1mm, y=1mm]
  \node[anchor=north west, inner sep=0pt] at (0,0) {\bksheet{3}{17}{%
1 & たて & よこ & \R{面積}{めん|せき} \\
2 & 3 & 4 & \SetCell{bg=bkiissel} \texttt{=A2*B2} \\
3 & 5 & 2 & \texttt{=A3*B3} \\
4 & 2 & 7 & \texttt{=A4*B4} \\
5 & 6 & 6 & \texttt{=A5*B5} \\}};
  \node[anchor=north west, inner sep=0pt] at (76,0) {\bksheet{3}{12}{%
1 & たて & よこ & \R{面積}{めん|せき} \\
2 & 3 & 4 & 12 \\
3 & 5 & 2 & 10 \\
4 & 2 & 7 & 14 \\
5 & 6 & 6 & 36 \\}};
  \node[anchor=south, font=\sffamily\footnotesize, text=bkGray] at (36,1.5) {\R{入}{はい}っている\R{式}{しき}};
  \node[anchor=south, font=\sffamily\footnotesize, text=bkGray] at (101,1.5) {\R{表示}{ひょう|じ}される\R{答}{こた}え};
\end{tikzpicture}
\bkfigcaption{\R{下}{した}へ1つコピーするごとに、\R{式}{しき}の\R{行}{ぎょう}の\R{番号}{ばん|ごう}が1ずつふえる。}\label{bk:2-18:copyfig}
\end{bkfigure}

コピーすると、\R{式}{しき}の\R{中}{なか}の\R{番地}{ばん|ち}も\R{同}{おな}じだけずれます。\R{下}{した}へ1つなら\R{行}{ぎょう}の\R{番号}{ばん|ごう}が1ふえ、\R{右}{みぎ}へ1つなら\R{列}{れつ}の\R{文字}{も|じ}が1つ\R{進}{すす}みます。

\textbf{なぜずれるの？}\quad \R{表計算}{ひょう|けい|さん}ソフトは、C2の \texttt{=A2*B2} を「A2とB2」という\R{名前}{な|まえ}ではなく、「\R{自分}{じ|ぶん}から\R{見}{み}て2つ\R{左}{ひだり}のセル×1つ\R{左}{ひだり}のセル」という\R{位置}{い|ち}の\R{関係}{かん|けい}で\R{覚}{おぼ}えています。\R{道案内}{みち|あん|ない}の「ここから\R{左}{ひだり}へ2\R{軒}{けん}\R{目}{め}」と\R{同}{おな}じで、\R{立}{た}っている\R{場所}{ば|しょ}が\R{変}{か}われば、\R{指}{さ}す\R{家}{いえ}も\R{変}{か}わります。

\begin{bkfigure}
\begin{tikzpicture}[x=1mm, y=1mm]
  \fill[bkiissel] (41,-4) rectangle (63,4);
  \fill[bkiissel] (41,-38) rectangle (63,-30);
  \foreach \r/\yy in {2/0, 3/-34} {
    \foreach \c/\xx/\t in {A/0/{A\r}, B/26/{B\r}, C/52/{C\r}} {
      \node[draw=bkLine, line width=0.6pt, minimum width=22mm, minimum height=8mm, font=\sffamily\small] (\c\r) at (\xx,\yy) {\t};}
    \node[anchor=west, font=\ttfamily\small, text=bkInk] at (66,\yy) {=A\r*B\r};}
  \foreach \r in {2,3} {
    \draw[bkBlue, line width=0.8pt, -{Stealth[length=3.5pt]}] ([yshift=1.5mm]C\r.north) to[out=150, in=30] node[pos=0.5, above, font=\sffamily\scriptsize, text=bkBlue!80!black, inner sep=1pt] {2つ\R{左}{ひだり}} ([yshift=1.5mm]A\r.north);
    \draw[bkOrange, line width=0.8pt, -{Stealth[length=3.5pt]}] ([yshift=-1.5mm]C\r.south) to[out=-150, in=-30] node[pos=0.5, below, font=\sffamily\scriptsize, text=bkOrange!80!black, inner sep=1pt] {1つ\R{左}{ひだり}} ([yshift=-1.5mm]B\r.south);}
\end{tikzpicture}
\bkfigcaption{C2もC3も、「\R{自分}{じ|ぶん}から2つ\R{左}{ひだり}×1つ\R{左}{ひだり}」という\R{同}{おな}じ\R{式}{しき}。}
\end{bkfigure}

だから、\R{式}{しき}をコピーするだけで、どの\R{行}{ぎょう}でも「\R{同}{おな}じ\R{行}{ぎょう}のたて×よこ」が\R{計算}{けい|さん}されます。\R{行}{ぎょう}が1000あっても、コピー1\R{回}{かい}ですみます。

\bksubsection{3}{ずらしたくない\R{番地}{ばん|ち}には \$ をつける}

\leavevmode\label{bk:2-18:dollar}ずれるのが\R{困}{こま}るときもあります。りんご1\R{個}{こ}の\R{値段}{ね|だん}120\R{円}{えん}をB1に\R{入}{い}れ、A\R{列}{れつ}の\R{個数}{こ|すう}から\R{代金}{だい|きん}を\R{出}{だ}します。B3に \texttt{=A3*B1} と\R{入}{い}れて\R{下}{した}へコピーすると、B4は \texttt{=A4*B2} になります。B2には「\R{代金}{だい|きん}」という\R{文字}{も|じ}しか\R{入}{はい}っていないので、\R{正}{ただ}しく\R{計算}{けい|さん}できません。

\R{値段}{ね|だん}のB1は、どの\R{行}{ぎょう}からでも\R{同}{おな}じセルを\R{使}{つか}いたいので、\texttt{\$B\$1} と\R{書}{か}きます。\$ は「ここは\R{動}{うご}かさない」という\R{印}{しるし}です。\R{列}{れつ}の\R{文字}{も|じ}の\R{前}{まえ}の \$ は\R{列}{れつ}を、\R{行}{ぎょう}の\R{番号}{ばん|ごう}の\R{前}{まえ}の \$ は\R{行}{ぎょう}を\R{止}{と}めます。\R{両方}{りょう|ほう}につけると、どの\R{向}{む}きにコピーしてもずれません。

\begin{bkfigure}
\begin{tikzpicture}[x=1mm, y=1mm]
  \node[anchor=north west, inner sep=0pt] at (0,0) {\bksheet{2}{24}{%
1 & 1\R{個}{こ}の\R{値段}{ね|だん} & \SetCell{bg=bkiissel} 120 \\
2 & \R{個数}{こ|すう} & \R{代金}{だい|きん} \\
3 & 1 & \texttt{=A3*\$B\$1} \\
4 & 2 & \texttt{=A4*\$B\$1} \\
5 & 3 & \texttt{=A5*\$B\$1} \\}};
  \node[anchor=north west, inner sep=0pt] at (76,0) {\bksheet{2}{18}{%
1 & 1\R{個}{こ}の\R{値段}{ね|だん} & 120 \\
2 & \R{個数}{こ|すう} & \R{代金}{だい|きん} \\
3 & 1 & 120 \\
4 & 2 & 240 \\
5 & 3 & 360 \\}};
\end{tikzpicture}
\bkfigcaption{A\R{列}{れつ}の\R{番地}{ばん|ち}はずれるが、\texttt{\$B\$1} はずれない。}
\end{bkfigure}

\R{値段}{ね|だん}が130\R{円}{えん}に\R{上}{あ}がったら、B1を130に\R{変}{か}えるだけで、\R{代金}{だい|きん}の\R{列}{れつ}が\R{全部}{ぜん|ぶ}\R{自動}{じ|どう}で\R{計算}{けい|さん}しなおされます。\R{変}{か}わるかもしれない\R{数}{かず}を1つのセルにまとめておくと、\R{便利}{べん|り}です。

\bksubsection{4}{\R{上}{うえ}の\R{行}{ぎょう}の\R{結果}{けっ|か}を\R{使}{つか}って、\R{下}{した}の\R{行}{ぎょう}を\R{計算}{けい|さん}する}

\leavevmode\label{bk:2-18:kurikaesi}\bkhatenaref の\R{貯金}{ちょ|きん}を\R{計算}{けい|さん}しましょう。きまりは「\R{今月}{こん|げつ}の\R{貯金}{ちょ|きん}は、\R{先月}{せん|げつ}の\R{貯金}{ちょ|きん}＋500」です。そこで、A\R{列}{れつ}に\R{月}{つき}、B\R{列}{れつ}に\R{貯金}{ちょ|きん}を\R{書}{か}く\R{表}{ひょう}を\R{作}{つく}ります。

\begin{enumerate}[label={\bkcircled{\arabic*}}, leftmargin=2.4\zw, itemsep=0pt, topsep=2pt]
\item 2\R{行目}{ぎょう|め}に、はじめ（0か\R{月}{げつ}）のようすを\R{入}{い}れる：A2に0、B2に3000。
\item 3\R{行目}{ぎょう|め}に、1つ\R{上}{うえ}の\R{行}{ぎょう}を\R{使}{つか}った\R{式}{しき}を\R{入}{い}れる：A3に \texttt{=A2+1}（\R{月}{つき}が1ふえる）、B3に \texttt{=B2+500}（\R{先月}{せん|げつ}＋500）。
\item A3とB3を\R{選}{えら}んで、\R{下}{した}へコピーする。
\end{enumerate}

\noindent\begin{minipage}[t]{0.56\linewidth}\vspace{0pt}\centering
\bksheet{2}{24}{%
1 & \R{月}{つき} & \R{貯金}{ちょ|きん}（\R{円}{えん}） \\
2 & 0 & 3000 \\
3 & \texttt{=A2+1} & \SetCell{bg=bkiissel} \texttt{=B2+500} \\
4 & \texttt{=A3+1} & \texttt{=B3+500} \\
5 & \texttt{=A4+1} & \texttt{=B4+500} \\}
\end{minipage}\hfill
\begin{minipage}[t]{0.42\linewidth}\vspace{0pt}\centering
\bksheet{2}{14}{%
1 & \R{月}{つき} & \R{貯金}{ちょ|きん} \\
2 & 0 & 3000 \\
3 & 1 & 3500 \\
4 & 2 & 4000 \\
$\vdots$ & $\vdots$ & $\vdots$ \\
16 & 14 & \SetCell{bg=bkiissel} 10000 \\
$\vdots$ & $\vdots$ & $\vdots$ \\
122 & 120 & 63000 \\}
\end{minipage}
\bktabcaption{3\R{行目}{ぎょう|め}の\R{式}{しき}を122\R{行目}{ぎょう|め}までコピーすると、1か\R{月}{げつ}ずつ\R{計算}{けい|さん}が\R{進}{すす}む。}

B3の \texttt{=B2+500} は、B2の3000を\R{使}{つか}って3500になります。B4にコピーされた \texttt{=B3+500} は、B3の\R{答}{こた}え3500を\R{使}{つか}って4000になります。こうして、\R{上}{うえ}から\R{順}{じゅん}に1\R{行}{ぎょう}ずつ、\R{前}{まえ}の\R{月}{つき}の\R{結果}{けっ|か}を\R{使}{つか}って\R{計算}{けい|さん}が\R{進}{すす}みます。リレーでバトンをわたしていくのと\R{同}{おな}じです。

\textbf{なぜ122\R{行目}{ぎょう|め}？}\quad 0か\R{月}{げつ}が2\R{行目}{ぎょう|め}にあるので、\R{行}{ぎょう}の\R{番号}{ばん|ごう}は、いつも「\R{月}{つき}の\R{数}{かず}＋2」です。14か\R{月}{げつ}は $14+2=16$ で16\R{行目}{ぎょう|め}、120か\R{月}{げつ}は $120+2=122$ で122\R{行目}{ぎょう|め}です。

\textbf{10\R{年}{ねん}は\R{何}{なん}か\R{月}{げつ}？}\quad 1\R{年}{ねん}あたり12か\R{月}{げつ}です。\R{単位}{たん|い}つきで\R{換算}{かん|さん}すると（2-10）、\R{年}{ねん}どうしが\R{約分}{やく|ぶん}されて、
\[ 10\,\cancel{\text{\R{年}{ねん}}}\times12\,\frac{\text{か\R{月}{げつ}}}{1\,\cancel{\text{\R{年}{ねん}}}}=120\,\text{か\R{月}{げつ}} \]
\bkhatenaref の\R{答}{こた}え：B\R{列}{れつ}を\R{下}{した}へたどると、10000が\R{出}{で}るのは16\R{行目}{ぎょう|め}、つまり\textbf{14か\R{月後}{げつ|ご}}です。120か\R{月後}{げつ|ご}（10\R{年後}{ねん|ご}）は122\R{行目}{ぎょう|め}で、\textbf{63000\R{円}{えん}}です。

\R{式}{しき}でも\R{確}{たし}かめられます。$x$ か\R{月後}{げつ|ご}の\R{貯金}{ちょ|きん}は $3000+500x$（\R{円}{えん}）です。10000\R{円}{えん}になるのは $3000+500x=10000$ のときで、2-16のように\R{解}{と}くと $500x=7000$、$x=14$。120か\R{月後}{げつ|ご}は $3000+500\times120=63000$ です。\R{表計算}{ひょう|けい|さん}と\R{同}{おな}じ\R{答}{こた}えになりました。

\R{表計算}{ひょう|けい|さん}では、\R{上}{うえ}の\R{行}{ぎょう}の\R{結果}{けっ|か}を\R{使}{つか}って\R{下}{した}の\R{行}{ぎょう}を\R{計算}{けい|さん}し、\R{式}{しき}をコピーしてくり\R{返}{かえ}せます。きまりを\R{変}{か}えれば、いろいろな\R{変}{か}わり\R{方}{かた}を\R{調}{しら}べられます。

\begin{bkcenter}
\begin{bktable}{colspec={Q[l,m,wd=50mm] Q[l,m,wd=22mm] Q[l,m,wd=50mm]}, rowsep=0.5pt, colsep=4pt, cells={font=\footnotesize}, row{1}={font=\sffamily\bfseries\footnotesize}}
\R{上}{うえ}の\R{行}{ぎょう}から\R{下}{した}の\R{行}{ぎょう}へのきまり & B3の\R{式}{しき} & B\R{列}{れつ}（B2から） \\
\R{毎月}{まい|つき}500\R{円}{えん}ずつふえる & \texttt{=B2+500} & 3000、3500、4000、4500、… \\
1\R{分}{ぷん}ごとに\qty{6}{L}ずつへる & \texttt{=B2-6} & 80、74、68、62、… \\
10\R{分}{ぷん}ごとに2\R{倍}{ばい}になる（2-7） & \texttt{=B2*2} & 100、200、400、800、… \\
\end{bktable}
\bktabcaption{きまりがちがっても、B3に\R{式}{しき}を1つ\R{書}{か}いてコピーするだけ。}
\end{bkcenter}

\bksubsection{5}{\R{表}{ひょう}からグラフを\R{作}{つく}る}

\leavevmode\label{bk:2-18:graph}\noindent\begin{minipage}[c]{0.47\linewidth}
\R{表}{ひょう}ができたら、グラフにしよう。
\begin{enumerate}[label={\bkcircled{\arabic*}}, leftmargin=2.2\zw, itemsep=0pt, topsep=2pt]
\item グラフにしたいセルのはんい（A1からB26など）を\R{選}{えら}ぶ。
\item グラフのメニューから、\R{点}{てん}を\R{線}{せん}でつなぐグラフを\R{選}{えら}ぶ。
\item \R{横}{よこ}の\R{軸}{じく}・たての\R{軸}{じく}に、\R{名前}{な|まえ}と\R{単位}{たん|い}をつける。
\end{enumerate}
\end{minipage}\hfill
\begin{minipage}[c]{0.5\linewidth}\centering
\begin{tikzpicture}
\begin{axis}[bkplot, width=64mm, height=42mm, xmin=0, xmax=26, ymin=0, ymax=17000, xtick={0,6,12,18,24}, ytick={0,5000,10000,15000},
  xlabel={\R{月}{つき}}, ylabel={\R{貯金}{ちょ|きん}（\R{円}{えん}）}]
  \addplot[bkcurve, line width=1pt, mark=*, mark size=1.1pt, samples at={0,1,...,24}] {3000+500*x};
\end{axis}
\end{tikzpicture}
\end{minipage}

\par\addvspace{4pt}
\R{毎月}{まい|つき}\R{同}{おな}じだけふえるので、\R{点}{てん}はまっすぐにならびます。2\R{倍}{ばい}ずつふえるきまりなら、\R{点}{てん}はだんだん\R{急}{きゅう}に\R{上}{あ}がっていきます。\R{表}{ひょう}の\R{数}{かず}を\R{変}{か}えると、グラフも\R{自動}{じ|どう}で\R{変}{か}わります。

\begin{taisetsu}
\R{式}{しき}は「＝」で\R{始}{はじ}め、セルの\R{番地}{ばん|ち}を\R{使}{つか}って\R{書}{か}く（×は *、÷は /）。\R{順番}{じゅん|ばん}は\R{算数}{さん|すう}と\R{同}{おな}じ\\
\R{式}{しき}をコピーすると\R{番地}{ばん|ち}もずれる。ずらしたくない\R{番地}{ばん|ち}は \texttt{\$B\$1}\\
\R{上}{うえ}の\R{行}{ぎょう}の\R{結果}{けっ|か}を\R{使}{つか}って\R{下}{した}の\R{行}{ぎょう}を\R{計算}{けい|さん}し、\R{式}{しき}をコピーしてくり\R{返}{かえ}せる
\end{taisetsu}

\begin{kotoba}
\small
\bkword{セル・\R{番地}{ばん|ち}}{\R{表計算}{ひょう|けい|さん}の\R{表}{ひょう}の1つ1つのますがセル。その\R{場所}{ば|しょ}を「B3」のように\R{表}{あらわ}したのが\R{番地}{ばん|ち}。}
\bkword{\R{式}{しき}のコピー}{\R{式}{しき}をほかのセルにうつすこと。\R{式}{しき}の\R{中}{なか}の\R{番地}{ばん|ち}も\R{同}{おな}じだけずれる。}
\end{kotoba}

\bkkigocards{\bkkigocard{\ttfamily *}{アスタリスク}{\R{表計算}{ひょう|けい|さん}やプログラムで、かけ\R{算}{ざん}×のかわりに\R{使}{つか}う\R{記号}{き|ごう}。}{asterisk}{2-18}}
{\bkkigocard{\ttfamily \$}{ドル}{\R{番地}{ばん|ち}の\R{列}{れつ}の\R{文字}{も|じ}・\R{行}{ぎょう}の\R{番号}{ばん|ごう}の\R{前}{まえ}につけると、コピーしてもずれない。}{dollar sign}{2-18}}

%% ---------------- 例題 ----------------
\bksection{reidai}{\R{例題}{れい|だい}でたしかめよう}

\begin{reidai}[\R{式}{しき}を\R{書}{か}く]
A2に\R{道}{みち}のり（m）、B2に\R{時間}{じ|かん}（s）が\R{入}{はい}っています。C2に\R{速}{はや}さを\R{出}{だ}す\R{式}{しき}を\R{書}{か}きましょう。A2が150、B2が30なら、C2には\R{何}{なに}が\R{出}{で}ますか。\par
\kangaekata \R{速}{はや}さ＝\R{道}{みち}のり÷\R{時間}{じ|かん}（2-3）。\R{道}{みち}のりはA2、\R{時間}{じ|かん}はB2なので、\R{番地}{ばん|ち}を\R{入}{い}れます。÷は / で\R{書}{か}きます。\par
\shiki \texttt{=A2/B2}。$150\div30=5$\par
\kotae \texttt{=A2/B2}、5（\R{単位}{たん|い}は $\mathrm{m}\div\mathrm{s}$ で\qty{5}{m/s}）
\end{reidai}

\begin{jibun}
A2にお\R{菓子}{か|し}1\R{個}{こ}の\R{重}{おも}さ（g）、B2に\R{個数}{こ|すう}が\R{入}{はい}っています。C2に\R{全体}{ぜん|たい}の\R{重}{おも}さを\R{出}{だ}す\R{式}{しき}を\R{書}{か}きましょう。A2が120、B2が8なら、C2には\R{何}{なに}が\R{出}{で}ますか。
\bkjans{\texttt{=A2*B2}、$120\times8=960$（\qty{960}{g}）}
\end{jibun}

\begin{reidai}[かっこが\R{要}{い}る\R{式}{しき}]
A2にケーキ1\R{個}{こ}の\R{値段}{ね|だん}300（\R{円}{えん}）、B2にジュース1\R{本}{ぽん}の\R{値段}{ね|だん}100（\R{円}{えん}）が\R{入}{はい}っています。ケーキとジュースを1つずつ、2\R{人分}{にん|ぶん}\R{買}{か}うときの\R{代金}{だい|きん}をC2に\R{出}{だ}す\R{式}{しき}を\R{書}{か}きましょう。\par
\kangaekata 1\R{人分}{にん|ぶん}は「ケーキ＋ジュース」で、それが2\R{人分}{にん|ぶん}です。「ケーキ＋ジュース」を1つのまとまりにしたいので、かっこでかこみます（2-11）。\par
\shiki \texttt{=(A2+B2)*2}。$(300+100)\times2=400\times2=800$\par
\hspace*{2.6\zw}かっこがないと \texttt{=A2+B2*2} で $300+100\times2=300+200=500$（ジュースだけが2\R{本}{ほん}になる）\par
\kotae \texttt{=(A2+B2)*2}、800（800\R{円}{えん}）
\end{reidai}

\begin{jibun}
A2に\R{長方形}{ちょう|ほう|けい}のたて（cm）、B2によこ（cm）が\R{入}{はい}っています。C2に\R{長方形}{ちょう|ほう|けい}のまわりの\R{長}{なが}さを\R{出}{だ}す\R{式}{しき}を\R{書}{か}きましょう。A2が5、B2が8なら、C2には\R{何}{なに}が\R{出}{で}ますか。
\bkjans{\texttt{=(A2+B2)*2}、$(5+8)\times2=26$（\qty{26}{cm}）}
\end{jibun}

\begin{reidai}[\R{式}{しき}のコピー]
D2に \texttt{=B2-C2} と\R{入}{い}れ、D6までコピーしました。D5の\R{式}{しき}は\R{何}{なん}ですか。\par
\kangaekata D2からD5へは\R{下}{した}へ3つ（$5-2=3$）なので、\R{式}{しき}の\R{中}{なか}の\R{行}{ぎょう}の\R{番号}{ばん|ごう}も3ふえます。\R{列}{れつ}の\R{文字}{も|じ}は\R{変}{か}わりません。\par
\shiki B2→B5、C2→C5\par
\kotae \texttt{=B5-C5}
\end{reidai}


\begin{reidai}[\$ をつけた\R{番地}{ばん|ち}（\R{時速}{じ|そく}を\R{秒速}{びょう|そく}に）]
A\R{列}{れつ}の3\R{行目}{ぎょう|め}から\R{下}{した}に\R{時速}{じ|そく}（km/h）が\R{入}{はい}っています。B3に\R{秒速}{びょう|そく}（m/s）を\R{出}{だ}す\R{式}{しき}を\R{書}{か}いて\R{下}{した}へコピーします。B1には\R{換算}{かん|さん}に\R{使}{つか}う\R{数}{かず}を\R{入}{い}れておきます。B1に\R{入}{い}れる\R{数}{かず}、B3の\R{式}{しき}、A5が72のときのB5の\R{値}{あたい}を\R{答}{こた}えましょう。\par
\kangaekata ① \R{時速}{じ|そく}\qty{72}{km}を\R{単位}{たん|い}つきで\R{秒速}{びょう|そく}にすると（2-10）、
\begin{align*}
\qty{72}{km/h}\times\qty{1000}{m/km}\div\qty{3600}{s/h}&=72\times\frac{1000}{3600}\,\mathrm{m/s}\\
&=\frac{72}{3.6}\,\mathrm{m/s}=\qty{20}{m/s}
\end{align*}
\hspace*{2.6\zw}$\dfrac{1000}{3600}=\dfrac{1}{3.6}$ なので、どの\R{時速}{じ|そく}も「÷3.6」で\R{秒速}{びょう|そく}になります。\par
\hspace*{2.6\zw}② 3.6は、どの\R{行}{ぎょう}でも\R{同}{おな}じ\R{数}{かず}なので、B1に1つだけ\R{入}{い}れ、\texttt{\$B\$1} で\R{使}{つか}います。\par
\shiki B1：3.6\quad B3：\texttt{=A3/\$B\$1}\par
\hspace*{2.6\zw}B5にコピーされた\R{式}{しき}は \texttt{=A5/\$B\$1} で、$72\div3.6=20$\par
\kotae B1は3.6、B3は \texttt{=A3/\$B\$1}、B5は20（\qty{20}{m/s}）
\end{reidai}

\begingroup\tcbset{equal height group=bkiisjA}%
\noindent\begin{minipage}[t]{0.49\linewidth}
\begin{jibun}
C3に \texttt{=A3*B3} と\R{入}{い}れ、C8までコピーしました。C7の\R{式}{しき}は\R{何}{なん}ですか。
\bkjans{C3からC7へは\R{下}{した}へ4つなので \texttt{=A7*B7}}
\end{jibun}
\end{minipage}\hfill
\begin{minipage}[t]{0.49\linewidth}
\begin{jibun}
B1に1 mの\R{値段}{ね|だん}150（\R{円}{えん}）、A\R{列}{れつ}の3\R{行目}{ぎょう|め}から\R{下}{した}に\R{長}{なが}さ（m）が\R{入}{はい}っています。B3に\R{代金}{だい|きん}の\R{式}{しき}を\R{書}{か}いて\R{下}{した}へコピーします。B3の\R{式}{しき}と、A4が2.5のときのB4の\R{値}{あたい}は？
\bkjans{\texttt{=A3*\$B\$1}、$2.5\times150=375$（375\nobreak\R{円}{えん}）}
\end{jibun}
\end{minipage}\par
\endgroup



\begin{reidai}[\R{上}{うえ}の\R{行}{ぎょう}の\R{結果}{けっ|か}を\R{使}{つか}う（ひき\R{算}{ざん}）]
\qty{80}{L}の\R{水}{みず}そうから、1\R{分}{ぷん}ごとに\qty{6}{L}ずつ\R{水}{みず}をぬきます。A2に0（\R{分}{ふん}）、B2に80を\R{入}{い}れ、B3の\R{式}{しき}を\R{下}{した}へコピーします。B3の\R{式}{しき}を\R{書}{か}き、\qty{50}{L}になるのは\R{何分後}{なん|ぷん|ご}かを\R{答}{こた}えましょう。\par
\kangaekata 1\R{分後}{ぷん|ご}の\R{水}{みず}は、1\R{分前}{ぷん|まえ}の\R{水}{みず}より\qty{6}{L}\R{少}{すく}ない。だから「1つ\R{上}{うえ}の\R{行}{ぎょう}－6」です。\par
\shiki B3：\texttt{=B2-6}\par
\hspace*{2.6\zw}B\R{列}{れつ}は 80、74、68、62、56、50（B2〜B7）\par
\hspace*{2.6\zw}50はB7。\R{行}{ぎょう}の\R{番号}{ばん|ごう}は「\R{分}{ふん}＋2」なので、$7-2=5$\par
\kotae \texttt{=B2-6}、5\R{分後}{ふん|ご}（\R{確}{たし}かめ：$80-6\times5=50$）
\end{reidai}

\begin{jibun}
\R{落}{お}ちるボールの\R{速}{はや}さは、0.1\R{秒}{びょう}ごとに $\qty{9.8}{m/s^2}\times\qty{0.1}{s}=\qty{0.98}{m/s}$ ずつふえます（2-10）。B2に0を\R{入}{い}れ、B3の\R{式}{しき}を\R{下}{した}へコピーします。B3の\R{式}{しき}と、0.5\R{秒後}{びょう|ご}（B7）の\R{速}{はや}さは？
\bkjans{\texttt{=B2+0.98}。0、0.98、1.96、2.94、3.92、4.9 で、\qty{4.9}{m/s}}
\end{jibun}

\begin{reidai}[2\R{倍}{ばい}ずつふえる（\R{時間}{じ|かん}の\R{換算}{かん|さん}つき）]
ある\R{細菌}{さい|きん}は、20\R{分}{ぷん}ごとに2\R{倍}{ばい}にふえます。B2にはじめの\R{数}{かず}50を\R{入}{い}れ、B3に「20\R{分後}{ぷん|ご}の\R{数}{かず}」の\R{式}{しき}を\R{書}{か}いて\R{下}{した}へコピーします。B3の\R{式}{しき}を\R{書}{か}き、2\R{時間後}{じ|かん|ご}の\R{数}{かず}が\R{何行目}{なん|ぎょう|め}に\R{出}{で}て、いくつになるかを\R{答}{こた}えましょう。\par
\kangaekata ① 1つ\R{下}{した}の\R{行}{ぎょう}は20\R{分後}{ぷん|ご}で、1つ\R{上}{うえ}の\R{行}{ぎょう}の2\R{倍}{ばい}。\par
\hspace*{2.6\zw}② 2\R{時間}{じ|かん}を\R{分}{ふん}に\R{換算}{かん|さん}して、20\R{分}{ぷん}が\R{何回}{なん|かい}あるかを\R{調}{しら}べます。
\[ \qty{2}{h}\times\qty{60}{min/h}=2\,\cancel{\mathrm{h}}\times\frac{60\,\mathrm{min}}{1\,\cancel{\mathrm{h}}}=\qty{120}{min},\qquad \qty{120}{min}\div\qty{20}{min}=6 \]
\hspace*{2.6\zw}20\R{分}{ぷん}が6\R{回}{かい}なので、2\R{倍}{ばい}が6\R{回}{かい}。B2から6\R{行}{ぎょう}\R{下}{した}のB8です。\par
\shiki B3：\texttt{=B2*2}\par
\hspace*{2.6\zw}B\R{列}{れつ}は 50、100、200、400、800、1600、3200（B2〜B8）\par
\kotae \texttt{=B2*2}、B8で3200（\R{確}{たし}かめ：$50\times2^6=50\times64=3200$、2-7）
\end{reidai}

\begin{jibun}
\R{厚}{あつ}さ\qty{0.1}{mm}の\R{紙}{かみ}は、1\R{回}{かい}\R{折}{お}るごとに\R{厚}{あつ}さが2\R{倍}{ばい}になります（2-7）。B2に0.1を\R{入}{い}れ、B3に「1\R{回}{かい}\R{折}{お}ったあと」の\R{式}{しき}を\R{書}{か}いて\R{下}{した}へコピーします。B3の\R{式}{しき}と、10\R{回}{かい}\R{折}{お}ったときの\R{厚}{あつ}さが\R{出}{で}るセル、その\R{厚}{あつ}さを\R{答}{こた}えましょう。
\bkjans{\texttt{=B2*2}。10\nobreak\R{回}{かい}\R{折}{お}るとB2から10\R{行}{ぎょう}\R{下}{した}のB12。$0.1\times2^{10}=102.4$ で\qty{102.4}{mm}}
\end{jibun}

\bkfukidashi{mina}{2\R{倍}{ばい}ずつふえると、10\R{回}{かい}で1000\R{倍}{ばい}\R{以上}{い|じょう}（$2^{10}=1024$）！\R{表計算}{ひょう|けい|さん}なら、こんな\R{大}{おお}きな\R{変化}{へん|か}もすぐに\R{確}{たし}かめられるね。}

\begin{reidai}[\R{差}{さ}の\R{割合}{わり|あい}で\R{変}{か}わる]
\qty{5}{\degreeCelsius}に\R{冷}{ひ}えた\R{飲}{の}み\R{物}{もの}を、\qty{25}{\degreeCelsius}の\R{部屋}{へや}に\R{置}{お}きます。1\R{分}{ぷん}ごとに「\R{部屋}{へや}と\R{飲}{の}み\R{物}{もの}の\R{温度}{おん|ど}の\R{差}{さ}」の\qty{10}{\percent}ずつ\R{温度}{おん|ど}が\R{上}{あ}がるとします。B2に5を\R{入}{い}れ、B3の\R{式}{しき}を\R{書}{か}いて\R{下}{した}へコピーします。B3の\R{式}{しき}と、B3・B4・B5の\R{値}{あたい}を\R{答}{こた}えましょう。\par
\kangaekata ① \R{差}{さ}は「\R{部屋}{へや}－\R{飲}{の}み\R{物}{もの}」で、\texttt{25-B2}。② その\qty{10}{\percent}は0.1をかけて \texttt{0.1*(25-B2)}（2-5）。\R{差}{さ}を\R{先}{さき}に\R{計算}{けい|さん}したいので、かっこが\R{要}{い}ります。③ それを\R{今}{いま}の\R{温度}{おん|ど}にたします。\par
\shiki B3：\texttt{=B2+0.1*(25-B2)}\par
\hspace*{2.6\zw}B3：$5+0.1\times(25-5)=5+0.1\times20=5+2=7$\par
\hspace*{2.6\zw}B4：$7+0.1\times(25-7)=7+1.8=8.8$\par
\hspace*{2.6\zw}B5：$8.8+0.1\times(25-8.8)=8.8+1.62=10.42$\par
\kotae \texttt{=B2+0.1*(25-B2)}。B3は7、B4は8.8、B5は10.42（℃）
\end{reidai}

\begin{jibun}
\qty{60}{\degreeCelsius}のお\R{湯}{ゆ}を、\qty{20}{\degreeCelsius}の\R{部屋}{へや}に\R{置}{お}きます。1\R{分}{ぷん}ごとに「お\R{湯}{ゆ}と\R{部屋}{へや}の\R{温度}{おん|ど}の\R{差}{さ}」の\qty{20}{\percent}ずつ\R{温度}{おん|ど}が\R{下}{さ}がるとします。B2に60を\R{入}{い}れたとき、B3の\R{式}{しき}と、B3・B4の\R{値}{あたい}を\R{答}{こた}えましょう。
\bkjans{\texttt{=B2-0.2*(B2-20)}。B3は $60-0.2\times40=52$、B4は $52-0.2\times32=45.6$}
\end{jibun}

\begin{bkfigure}
\begin{tikzpicture}
\begin{axis}[bkplot, width=86mm, height=44mm, xmin=0, xmax=31, ymin=0, ymax=28, xtick={0,10,20,30}, ytick={5,10,15,20,25},
  xlabel={\R{分}{ふん}}, ylabel={\R{温度}{おん|ど}（℃）}]
  \addplot[bkthin, dashed, domain=0:30] {25};
  \addplot[bkcurve, only marks, mark=*, mark size=1pt, samples at={0,1,...,30}] {25-20*0.9^x};
  \node[anchor=north west, font=\sffamily\footnotesize, text=bkGray, fill=white, inner sep=1pt] at (axis cs:15,23.5) {\R{部屋}{へや}の\R{温度}{おん|ど}（\qty{25}{\degreeCelsius}）};
\end{axis}
\end{tikzpicture}
\bkfigcaption{\R{例題}{れい|だい}7の\R{式}{しき}を31\R{行}{ぎょう}コピーしてグラフにしたもの。\R{差}{さ}が\R{小}{ちい}さくなるほど\R{上}{あ}がり\R{方}{かた}がゆっくりになり、\R{部屋}{へや}の\R{温度}{おん|ど}に\R{近}{ちか}づく。}
\end{bkfigure}

\begin{reidai}[\R{合計}{ごう|けい}と\R{平均}{へい|きん}を\R{出}{だ}す\R{命令}{めい|れい}]
B2からB6に、5\R{回}{かい}のテストの\R{点}{てん}（72、85、64、90、79）が\R{入}{はい}っています。B7に\R{合計}{ごう|けい}、B8に\R{平均}{へい|きん}を\R{出}{だ}す\R{式}{しき}を\R{書}{か}きましょう。\par
\noindent\begin{minipage}[t]{0.58\linewidth}
\kangaekata \texttt{=B2+B3+B4+B5+B6} でもできますが、\R{行}{ぎょう}が\R{多}{おお}いと\R{大変}{たい|へん}です。\R{表計算}{ひょう|けい|さん}には、はんいの\R{数}{かず}をまとめてたす \texttt{SUM}（サム、\R{合計}{ごう|けい}）という\R{命令}{めい|れい}があります。\texttt{B2:B6} は「B2からB6まで」という\R{意味}{い|み}です。\R{平均}{へい|きん}は\R{合計}{ごう|けい}÷\R{個数}{こ|すう}（2-6）なので、\texttt{=B7/5} か、\R{平均}{へい|きん}の\R{命令}{めい|れい} \texttt{AVERAGE}（アベレージ）を\R{使}{つか}います。\par
\shiki B7：\texttt{=SUM(B2:B6)}\quad $72+85+64+90+79=390$\par
\hspace*{2.6\zw}B8：\texttt{=AVERAGE(B2:B6)}\quad $390\div5=78$\par
\kotae B7は \texttt{=SUM(B2:B6)} で390、B8は \texttt{=AVERAGE(B2:B6)}（または \texttt{=B7/5}）で78
\end{minipage}\hfill
\begin{minipage}[t]{0.38\linewidth}\vspace{2pt}\raggedleft
\bksheet{1}{35}{%
1 & \R{点}{てん} \\
2 & 72 \\
3 & 85 \\
4 & 64 \\
5 & 90 \\
6 & 79 \\
7 & \SetCell{bg=bkiissel}{\footnotesize\texttt{=SUM(B2:B6)}} \\
8 & \SetCell{bg=bkiissel}{\footnotesize\texttt{=AVERAGE(B2:B6)}} \\}
\end{minipage}
\end{reidai}

\begin{jibun}
A2からA5に、4\R{日間}{か|かん}の\R{読書}{どく|しょ}の\R{時間}{じ|かん}（12、15、9、14 \R{分}{ふん}）が\R{入}{はい}っています。A6に\R{合計}{ごう|けい}、A7に\R{平均}{へい|きん}を\R{出}{だ}す\R{式}{しき}と、その\R{値}{あたい}を\R{答}{こた}えましょう。
\bkjans{A6は \texttt{=SUM(A2:A5)} で50、A7は \texttt{=AVERAGE(A2:A5)}（または \texttt{=A6/4}）で12.5}
\end{jibun}

\bkfukidashi{sora}{\texttt{SUM} や \texttt{AVERAGE} のような\R{命令}{めい|れい}を\R{使}{つか}うと、1000\R{行}{ぎょう}の\R{合計}{ごう|けい}も\R{平均}{へい|きん}も1つの\R{式}{しき}で\R{出}{だ}せるんだ。}

%% ---------------- まちがえやすいところ ----------------
\bksection{tsumazuki}{まちがえやすいところをチェック}

\begin{tsumazuki}{「＝」をつけわすれる}
\bkng{C2に \texttt{A2*B2} と\R{入}{い}れる（\R{文字}{も|じ}のまま\R{表示}{ひょう|じ}される）}
\bkok{C2に \texttt{=A2*B2} と\R{入}{い}れる（\R{計算}{けい|さん}した\R{答}{こた}えが\R{出}{で}る）}
「＝」は、「ここに\R{計算}{けい|さん}の\R{答}{こた}えを\R{入}{い}れて」という\R{合図}{あい|ず}です。
\end{tsumazuki}

\begin{tsumazuki}{\$ をつけわすれて、コピーする}
\bkng{B3に \texttt{=A3*B1} → B4にコピーすると \texttt{=A4*B2}（B2は「\R{代金}{だい|きん}」の\R{文字}{も|じ}）}
\bkok{B3に \texttt{=A3*\$B\$1} → B4にコピーすると \texttt{=A4*\$B\$1}}
\R{同}{おな}じセルを\R{使}{つか}い\R{続}{つづ}けたい\R{番地}{ばん|ち}には、\$ をつけます。
\end{tsumazuki}

\begin{tsumazuki}{かっこをつけわすれる}
\bkng{B3に \texttt{=B2-0.05*B2-20}（80なら $80-4-20=56$）}
\bkok{B3に \texttt{=B2-0.05*(B2-20)}（80なら $80-0.05\times60=77$）}
かっこがないと、*を\R{先}{さき}に\R{計算}{けい|さん}し（2-11）、20を\R{最後}{さい|ご}にひいてしまいます。「\R{差}{さ}の5\%」の「\R{差}{さ}」を1つのまとまりにするには、かっこが\R{要}{い}ります。
\end{tsumazuki}

\begin{tsumazuki}{\R{自分}{じ|ぶん}のセルを\R{式}{しき}に\R{入}{い}れてしまう}
\bkng{B3に \texttt{=B3+500}（「\R{循環参照}{じゅん|かん|さん|しょう}」の\R{注意}{ちゅう|い}が\R{出}{で}る）}
\bkok{B3に \texttt{=B2+500}（1つ\R{上}{うえ}の\R{行}{ぎょう}を\R{使}{つか}う）}
B3の\R{答}{こた}えを\R{出}{だ}すのにB3の\R{答}{こた}えが\R{要}{い}ると、いつまでも\R{計算}{けい|さん}が\R{終}{お}わりません。「\R{先月}{せん|げつ}＋500」の「\R{先月}{せん|げつ}」は、1つ\R{上}{うえ}の\R{行}{ぎょう}です。
\end{tsumazuki}


\begin{tsumazuki}{\R{番地}{ばん|ち}の\R{列}{れつ}と\R{行}{ぎょう}を\R{逆}{ぎゃく}に\R{書}{か}く}
\bkng{B\R{列}{れつ}の3\R{行目}{ぎょう|め}を「3B」と\R{書}{か}く}
\bkok{B\R{列}{れつ}の3\R{行目}{ぎょう|め}は「B3」（\R{列}{れつ}の\R{文字}{も|じ}が\R{先}{さき}、\R{行}{ぎょう}の\R{番号}{ばん|ごう}があと）}
「3B」と\R{書}{か}くと、\R{表計算}{ひょう|けい|さん}ソフトは\R{番地}{ばん|ち}だとわかりません。
\end{tsumazuki}



\clearpage
%% ---------------- 物理・工学 ----------------
\begin{butsuri}
\textbf{お\R{湯}{ゆ}の\R{冷}{さ}め\R{方}{かた}を\R{予想}{よ|そう}する}\par
\R{熱}{あつ}いお\R{湯}{ゆ}は、\R{部屋}{へや}の\R{温度}{おん|ど}との\R{差}{さ}が\R{大}{おお}きいほど\R{速}{はや}く\R{冷}{さ}めます。\R{差}{さ}が\R{小}{ちい}さくなると、\R{冷}{さ}め\R{方}{かた}もゆっくりになります。あるカップでは、1\R{分}{ぷん}ごとに「お\R{湯}{ゆ}と\R{部屋}{へや}（\qty{20}{\degreeCelsius}）の\R{温度}{おん|ど}の\R{差}{さ}」の\qty{5}{\percent}ずつ\R{下}{さ}がりました。

\R{式}{しき}を1つずつ\R{組}{く}み\R{立}{た}てます。① \R{差}{さ}は「お\R{湯}{ゆ}－\R{部屋}{へや}」で \texttt{B2-20}。② その\qty{5}{\percent}は \texttt{0.05*(B2-20)}。③ それを\R{今}{いま}の\R{温度}{おん|ど}からひく。B2に80を\R{入}{い}れ、B3に \texttt{=B2-0.05*(B2-20)} と\R{入}{い}れて\R{下}{した}へコピーします。\R{手}{て}で\R{確}{たし}かめると、B3は $80-0.05\times60=80-3=77$、B4は $77-0.05\times57=77-2.85=74.15$ です。

\noindent\begin{minipage}[c]{0.4\linewidth}\centering
\bksheet{2}{15}{%
1 & \R{分}{ふん} & ℃ \\
2 & 0 & 80 \\
3 & 1 & 77 \\
4 & 2 & 74.15 \\
5 & 3 & 71.44 \\
$\vdots$ & $\vdots$ & $\vdots$ \\
12 & 10 & 55.92 \\}
\par\vspace{2pt}{\footnotesize\sffamily（\R{小数}{しょう|すう}\R{第}{だい}3\R{位}{い}を\R{四捨五入}{し|しゃ|ご|にゅう}）}
\end{minipage}\hfill
\begin{minipage}[c]{0.58\linewidth}\centering
\begin{tikzpicture}
\begin{axis}[bkplot, width=74mm, height=46mm, xmin=0, xmax=33, ymin=0, ymax=92, xtick={0,10,20,30}, ytick={20,40,60,80},
  xlabel={\R{分}{ふん}}, ylabel={\R{温度}{おん|ど}（℃）}]
  \addplot[bkthin, dashed, domain=0:31] {20};
  \addplot[bkcurve2, only marks, mark=*, mark size=1pt, samples at={0,1,...,30}] {20+60*0.95^x};
  \node[anchor=north west, font=\sffamily\footnotesize, text=bkGray, fill=white, inner sep=1pt] at (axis cs:12,17) {\R{部屋}{へや}の\R{温度}{おん|ど}};
\end{axis}
\end{tikzpicture}
\end{minipage}

\par\addvspace{4pt}
はじめは1\R{分}{ぷん}に\qty{3}{\degreeCelsius}も\R{下}{さ}がりますが、10\R{分後}{ぷん|ご}には1\R{分}{ぷん}に\R{約}{やく}\qty{1.8}{\degreeCelsius}しか\R{下}{さ}がらず、だんだんゆっくり\R{冷}{さ}めて、\R{部屋}{へや}の\R{温度}{おん|ど}に\R{近}{ちか}づきます。「\R{今}{いま}の\R{状態}{じょう|たい}から、\R{少}{すこ}し\R{先}{さき}の\R{状態}{じょう|たい}を\R{計算}{けい|さん}する」ことをくり\R{返}{かえ}すこの\R{方法}{ほう|ほう}は、\R{第}{だい}5\R{巻}{かん}（5-8）・\R{第}{だい}6\R{巻}{かん}（6-2 オイラー\R{法}{ほう}）で\R{本格的}{ほん|かく|てき}に\R{使}{つか}います。ロケットの\R{進}{すす}み\R{方}{かた}や\R{天気予報}{てん|き|よ|ほう}も、コンピュータがこの\R{方法}{ほう|ほう}で\R{計算}{けい|さん}しています。
\end{butsuri}




\clearpage
%% ---------------- 章末問題 ----------------
\bkendproblems
\bkproblemgroup{\R{基本}{き|ほん}}{2}

\bkproblem \R{右}{みぎ}の\R{表}{ひょう}を\R{見}{み}て\R{答}{こた}えましょう。\par
\noindent\begin{minipage}[c]{0.5\linewidth}
（1）\nobreak B3に\R{入}{はい}っている\R{数}{かず}は\R{何}{なん}ですか。\par
（2）\nobreak 25が\R{入}{はい}っているセルの\R{番地}{ばん|ち}は\R{何}{なん}ですか。
\end{minipage}\hfill
\begin{minipage}[c]{0.46\linewidth}\centering
\bksheet{3}{12}{%
1 & 4 & 9 & 16 \\
2 & 7 & 25 & 3 \\
3 & 11 & 8 & 30 \\}
\end{minipage}
\bkpans{（1）\nobreak 8\quad（2）\nobreak B2}

\bkproblem A2とB2の\R{数}{かず}を\R{使}{つか}う\R{式}{しき}を、C2に\R{入}{い}れる\R{形}{かたち}で\R{書}{か}きましょう。\par
（1）\nobreak \R{和}{わ}（たし\R{算}{ざん}の\R{答}{こた}え）\quad（2）\nobreak \R{積}{せき}（かけ\R{算}{ざん}の\R{答}{こた}え）\quad（3）\nobreak A2をB2でわった\R{商}{しょう}
\bkpans{（1）\nobreak \texttt{=A2+B2}\quad（2）\nobreak \texttt{=A2*B2}\quad（3）\nobreak \texttt{=A2/B2}}

\bkproblem A2に10、B2に6が\R{入}{はい}っています。C2に \texttt{=A2+B2/2} と\R{入}{い}れると、\R{何}{なに}が\R{出}{で}ますか。また、A2とB2の\R{平均}{へい|きん}を\R{出}{だ}したいとき、\R{式}{しき}をどう\R{直}{なお}せばよいですか。
\bkpans{$10+6\div2=13$。\R{平均}{へい|きん}は \texttt{=(A2+B2)/2} で8}

\bkproblem D2に \texttt{=B2*C2} と\R{入}{い}れ、D6までコピーしました。D4の\R{式}{しき}は\R{何}{なん}ですか。
\bkpans{\texttt{=B4*C4}}

\bkproblem C3に \texttt{=B3*\$E\$1} と\R{入}{い}れ、C7までコピーしました。C6の\R{式}{しき}は\R{何}{なん}ですか。
\bkpans{\texttt{=B6*\$E\$1}}

\bkproblem B2に2000を\R{入}{い}れ、B3に \texttt{=B2+150} と\R{入}{い}れて、\R{下}{した}へコピーしました。B6に\R{表示}{ひょう|じ}される\R{数}{かず}は\R{何}{なん}ですか。
\bkpans{B3が2150、B4が2300、B5が2450、B6が2600}

\bkproblem \begin{minipage}[t]{0.4\linewidth}
C2に \texttt{=A2*B2} と\R{入}{い}れ、C4までコピーしました。C3とC4に\R{表示}{ひょう|じ}される\R{数}{かず}は\R{何}{なん}ですか。
\end{minipage}\hfill
\begin{minipage}[t]{0.52\linewidth}\vspace{0pt}\raggedleft
\bksheet{3}{16}{%
1 & たて（cm） & よこ（cm） & \R{面積}{めん|せき}（cm²） \\
2 & 4 & 6 & \texttt{=A2*B2} \\
3 & 7 & 3 & \\
4 & 2.5 & 8 & \\}
\end{minipage}\par
\bkpans{C3は \texttt{=A3*B3} で21、C4は \texttt{=A4*B4} で20}

\bkproblemgroup{チャレンジ}{2}

\bkproblem ある\R{細菌}{さい|きん}は、10\R{分}{ぷん}ごとに2\R{倍}{ばい}にふえます。B2にはじめの\R{数}{かず}100を\R{入}{い}れ、B3に\R{式}{しき}を\R{書}{か}いて\R{下}{した}へコピーします。B3の\R{式}{しき}と、1\R{時間後}{じ|かん|ご}の\R{数}{かず}を\R{答}{こた}えましょう（\R{換算}{かん|さん}は\R{単位}{たん|い}つきで）。
\bkpans{\texttt{=B2*2}。$\qty{1}{h}\times\qty{60}{min/h}=\qty{60}{min}$、$\qty{60}{min}\div\qty{10}{min}=6$ で2\R{倍}{ばい}が6\nobreak\R{回}{かい}。B8で $100\times2^6=6400$}

\bkproblem ゆうさんは、B3に \texttt{=A3*B1} と\R{入}{い}れて\R{下}{した}へコピーしたところ、B4から\R{下}{した}の\R{答}{こた}えがおかしくなりました。なぜですか。B3の\R{式}{しき}をどう\R{直}{なお}せばよいですか。
\bkpans{コピーするとB1がB2・B3…とずれるから。\texttt{=A3*\$B\$1} に\R{直}{なお}す}

\bkproblem お\R{湯}{ゆ}は、1\R{分}{ぷん}ごとに「\R{部屋}{へや}（\qty{20}{\degreeCelsius}）との\R{差}{さ}」の\qty{10}{\percent}ずつ\R{下}{さ}がります。B2に80、B3に \texttt{=B2-0.1*(B2-20)} と\R{入}{い}れて\R{下}{した}へコピーしました。B3・B4・B5の\R{数}{かず}を\R{求}{もと}めましょう。
\bkpans{B3 $80-0.1\times60=74$、B4 $74-0.1\times54=68.6$、B5 $68.6-0.1\times48.6=63.74$}

\begin{furikaeri}
\begin{checklist}
\item 「＝」の\R{式}{しき}を\R{書}{か}き、コピーしてくり\R{返}{かえ}せる。\hfill{\footnotesize\sffamily\color{bkTheme}→\pageref{bk:2-18:shiki}ページ}
\item ずらしたくない\R{番地}{ばん|ち}に \$ をつけて\R{使}{つか}える。\hfill{\footnotesize\sffamily\color{bkTheme}→\pageref{bk:2-18:dollar}ページ}
\item \R{上}{うえ}の\R{行}{ぎょう}の\R{結果}{けっ|か}で\R{下}{した}の\R{行}{ぎょう}を\R{計算}{けい|さん}し、グラフにできる。\hfill{\footnotesize\sffamily\color{bkTheme}→\pageref{bk:2-18:kurikaesi}ページ}
\end{checklist}
\end{furikaeri}

\bkfukidashi{sora}{コピー1\R{回}{かい}で、1000\R{回}{かい}の\R{計算}{けい|さん}ができたね。\R{次}{つぎ}の2-19では、\R{計算}{けい|さん}の\R{手順}{て|じゅん}を\R{自分}{じ|ぶん}で\R{書}{か}く「プログラム」を\R{学}{まな}ぶよ。}
